\honest{The Bottom Line}{Eliezer Yudkowsky}{2007}

Two sealed boxes are up for auction, and one of them holds a diamond. There are clues, none of them certain. A blue stamp is more common on boxes with diamonds, and I suspect, without being sure, that shiny boxes are always empty.

A clever arguer offers to argue for whichever owner pays more, and box B's owner wins. The arguer first writes at the bottom of a sheet of paper, ``And therefore, box B contains the diamond!'' Then the arguer fills the lines above with the clues that favor B, and leaves out the clues that favor A.

The moment the conclusion was written, its relation to the boxes was fixed. To help you see this, I bring in Everett branches and Tegmark duplicates. \nb{Ordinary probability would do: over many auctions like this one, how often is the diamond in B when the arguer's page says B? The physics adds nothing.}

I warn that a literate reader might think the words ``box B contains the diamond'' mean that box B contains the diamond. My evidence is the Stroop test, in which people shown the word ``green'' printed in red ink often say ``green.'' \nb{That test shows that reading is hard to switch off. It does not show that people believe what they read.} ``It helps to be illiterate.''

The key step is this. To us, what a thing tells us is its entanglement with other things. Whatever is written on the lines above the conclusion, it cannot change what the conclusion tells you, and the arguer's handwriting is evidence ``only of which owner paid the higher bid.'' I do not say what the listed clues are worth. \nb{They are still evidence: box B does have a blue stamp. What the arguer hides is the clues for A, so a reader should discount the list, not ignore it. Yudkowsky made this point the next day, in ``What Evidence Filtered Evidence?''}

Along the way I raise a question: what the arguer's conclusion tells you depends on which owners tend to hire arguers. Perhaps owners who think they have the better case hire advertisers more often; perhaps owners who fear their box is worse bid higher. If the owners do not understand the clues, the page tells you about their finances and bidding habits, and nothing about the diamond. I leave the rest ``to you.''

Then a curious person lists the clues for both boxes and writes ``85\%'' at the bottom. That page does depend on the clues: in worlds where the clues differ, a different number is written at the bottom. The two pages may sound alike when read aloud, but they are evidence of different things.

Near the end I make the point of the essay: you may be the clever arguer yourself. If your brakes squeal and you do not want to pay for a repair, you may look for reasons why the car is fine. What decides your fate is the rule that chose which conclusion to argue for, here ``Never repair anything expensive.'' I define your effectiveness as the share of you that survives across Everett branches. I do not tell you how to find out which rule is writing your own conclusions. \nb{From the inside, motivated reasoning feels like honest reasoning.}

Last, a caution: do not use this essay to dismiss your opponents as clever arguers. Saying ``My opponent is a clever arguer'' is itself a clever argument, if what you want is to keep the beliefs you started with. ``The world's cleverest arguer may point out that the Sun is shining, and yet it is still probably daytime.'' \nb{That is true. It concedes the point of the earlier note: what the arguer lists can still be evidence.}
